\chapter*{Abstract}
\setheader{Abstract}

ERP-Systeme können für kleine und mittlere Unternehmen eine wichtige technologische Grundlage der digitalen 
Transformation darstellen, da sie Geschäftsprozesse integrieren und eine gemeinsame Datenbasis schaffen. Ihre 
Einführung führt jedoch nicht automatisch zu einer transformierenden Wirkung. Ziel dieser Arbeit ist es daher, 
zu untersuchen, welche Voraussetzungen KMU erfüllen müssen, damit ERP-Systeme als technologische Basis digitaler 
Transformation wirksam werden können.

Auf Grundlage einer literaturbasierten Analyse werden zentrale Erfolgsbedingungen für die ERP-Wirksamkeit in KMU 
identifiziert. Dazu zählen ERP-fähige Prozesse, integrierte und nutzbare Daten, realistisch verfügbare Ressourcen, 
Kompetenzen und Akzeptanz der Mitarbeitenden sowie die strategische Einbettung des ERP-Systems. Diese Voraussetzungen 
werden im entwickelten ERP-KMU-Fit-Modell als Prozess-, Daten-, Ressourcen-, Kompetenz- und Strategie-Dimension zusammengeführt.

Die Arbeit zeigt, dass ERP-Systeme insbesondere dann transformativ wirken können, wenn zwischen den Anforderungen des 
Systems und den organisatorischen, technologischen sowie strategischen Voraussetzungen eines KMU eine ausreichende Passung 
besteht. Das Modell dient dabei als konzeptionelles Orientierungsinstrument, um mögliche Handlungsbedarfe vor oder während 
einer ERP-Einführung systematisch sichtbar zu machen. Es ersetzt jedoch keine empirische Erfolgsvorhersage, da 
branchenspezifische Unterschiede, Unternehmenskultur und weitere kontextuelle Faktoren nur begrenzt abgebildet werden können.\\  

\medskip
\noindent\textbf{Schlüsselwörter:} \emph{ERP-Systeme, KMU, digitale Transformation, ERP-Einführung, Erfolgsfaktoren, Fit-Modell}

